\documentclass[11pt]{article}
\usepackage{fontspec}
\usepackage[french]{babel}
\usepackage{amsmath}
\usepackage{amssymb}
\usepackage{amsthm}
\usepackage{mathrsfs}
\usepackage{tikz-cd}
\usepackage{enumitem}

\newcommand{\printedpage}[1]{\marginpar{\footnotesize p.~#1}}
\newcommand{\msp}[1]{\marginpar{\footnotesize [#1]}}
\newcommand{\circled}[1]{\tikz[baseline=(c.base)]\node[draw,circle,inner sep=1pt](c){#1};}

\begin{document}

\printedpage{54}
% head: A. Grothendieck : Dérivateurs, Ch. XVI, Édition des Universités de Montpellier II et Paris VII
\msp{page 114}
(\textit{cf.} p. 110, cas b)) et \textit{de plus} Der (3). (Donc dans le cas d'un triple de Quillen \textit{très} spécial, le prédérivateur obtenu est bel et bien un dérivateur.) Je vais expliciter ici la démonstration.

Je vais désigner par $(\mathcal{M},\Sigma)$ la catégorie de modèles donnée ($\Sigma\subset\operatorname{Fl}\mathcal{M}$), par $\mathbf{D}=\mathbf{D}_{\mathcal{M},W}$ le prédérivateur correspondant, et par
\[ W=W_D \] % [sic] W_D -> W_{\mathbf{D}}, index printed as italic D, elsewhere bold D
le localiseur dans Cat associé, formé des ``$\mathbf{D}$-équivalences''. On voit tout de suite que $W$ est fortement saturé (vrai pour tout $W_{\mathbf{D}}$, $\mathbf{D}$ un prédérivateur), \textit{i.e.} satisfait Loc (1 ter), a fortiori Loc (1), d'autre part il satisfait à Loc (3 bis$^*$), comme conséquence du fait que $\mathbf{D}$ satisfait à Loc (4$^*$, 5$^*$). A fortiori il satisfait à Loc (3$^*$), et pour voir que tout $X$ contractile dans Cat est $W$-asphérique, \textit{i.e.} $X\longrightarrow e$ dans $W$, on sait (\textit{cf.} §1)
\msp{page 115}
qu'il suffit de prouver qu'il en est ainsi si $X$ a un objet initial (axiome Loc (2$^*$) dual de l'axiome Loc (2)). À ce moment, on saura que $W$ est un localiseur fondamental, et de plus il satisfait Loc (3$^*$ bis), donc aussi Loc (3 bis) (forme forte de l'axiome de localisation).

Donc \textit{on peut supposer} que si $S$ a un objet initial, alors $S$ est $W$-asphérique, \textit{i.e.} $S\longrightarrow e$ est dans $W$, \textit{i.e.} $p^*:\mathbf{D}(e)\longrightarrow\mathbf{D}(S)$ est pleinement fidèle. Il revient au même de dire que la flèche d'adjonction
\[ p_!p^*\xi\longrightarrow\xi \]
est un isomorphisme, pour tout $\xi$ dans $\mathbf{D}(e)$. Or le calcul de $p_!(\eta)$ (pour $\eta\in\operatorname{Ob}\mathbf{D}(S)$, \textit{i.e.} $\eta\in\operatorname{Ob}\mathcal{M}(S)$) se fait ainsi : on choisit une flèche
\[ \eta'\xrightarrow{\Sigma_S}\eta \]
avec $\eta'$ ``cofibrant'' (pour la structure à cofibrations pertinente sur $\mathcal{M}(S)$), et on
\msp{page 116}
pose \footnote[54]{opération $p_!$ au niveau de $\mathcal{M}(S)$.}
\[ p_!(\operatorname{ho}_{W,S}(\eta))=\operatorname{ho}_W(p_!(\eta')). \]
Quand $\eta=p^*(\xi)$, la flèche d'adjonction pertinente dans $\operatorname{Ho}_\Sigma=\mathbf{D}(e)$ s'obtient par composition
\[ p_!(\eta')\longrightarrow p_!(p^*(\xi))\xrightarrow{\text{adj.}}\xi, \tag*{$(*)$} \]
où ici $p_!$ est encore pris dans le sens ordinaire ($\mathcal{M}(S)\longrightarrow\mathcal{M}(e)=\mathcal{M}$). Or $S$ ayant un objet initial 0, on a
\[ \begin{aligned}
p_!(\eta') &\ \simeq\ \eta'(0)\\
p_!(p^*(\xi)) &\ \simeq\ p^*(\xi)(0)\ \simeq\ \xi,
\end{aligned} \]
et la flèche d'adjonction dans $(*)$ est l'identité, la première flèche est la flèche $\eta'(0)\longrightarrow\eta(0)$ induite par $\eta'\longrightarrow\eta$, elle est dans $\Sigma$, donc $(*)$ est dans $\Sigma$, donc c'est un isomorphisme dans $\operatorname{Ho}_\Sigma=\mathbf{D}(e)$, q.e.d.

Le même argument marche d'ailleurs si $\mathcal{M}$ est seulement une catégorie à cofibrations ou à fibrations, mais en se bornant alors au cas où $S$ est un ensemble ordonné fini.
% foot: Version 18 mars 2019

% Original hard page break suppressed for the reflowed reader.
\printedpage{55}
% head: A. Grothendieck : Dérivateurs, Ch. XVI, Édition des Universités de Montpellier II et Paris VII
\msp{page 117}
Je vais dégager à part ce qu'on vient d'obtenir :

\textbf{Proposition 7.7.1.} \textit{Soit} $(\mathcal{M},W)$ \textit{une catégorie de modèles, associée à un triple de Quillen clos} spécial $(W,C,F)$\textit{. Alors le prédérivateur} $\mathbf{D}=\mathbf{D}_{\mathcal{M},W}$ \textit{sur} Cat \textit{satisfait aux axiomes} Der (1, 2, 3, 4$^*$, 5$^*$, 6$^*$)\textit{. Si de plus le triple de Quillen est} très \textit{spécial} (\textit{donc en particulier} $C=\operatorname{Mono}(\mathcal{M})$), \textit{alors il satisfait de plus à} Der (4, 5, 6), i.e. \textit{c'est un dérivateur au plein sens du terme. D'autre part, le localiseur} $W=W_{\mathbf{D}}$ \textit{dans} Cat \textit{associé à} $\mathbf{D}$ (\textit{dans le cas d'un triple de Quillen spécial sans plus}) \textit{est un} localiseur fondamental, \textit{satisfaisant de plus} Loc (1 ter, 3 bis).

\textbf{N.B.} Ce qui est dit ici pour le cas d'un triple de Quillen très spécial (qui ne fait pas appel à la théorie des localiseurs
\msp{page 118}
fondamentaux dans Cat) montre à nouveau que si $W$ localiseur fondamental dans Cat satisfait à Loc (6$'$ a, 8, 9), alors $\mathrm{HOT}_W$ est un vrai dérivateur. On a vu de plus dans 7.6 (page 111) que l'on a alors $W=W_{\mathrm{HOT}_W}$ \footnote[55]{Mais c'était sous l'hypothèse Loc (3 bis) ! Gare au cercle vicieux.}, donc il s'ensuit :

\textbf{Corollaire 7.7.2 :} \textit{On a l'implication}
\[ \text{Loc } (6'\text{ a, 8, 9})\ \Longrightarrow\ \text{Loc(1 ter, 3 bis)}\ \text{\footnotemark[56]}. \]
\footnotetext[56]{Loc (3 bis) pas prouvé, probablement faux.}

Notons aussi, pour mémoire :

\textbf{Corollaire 7.7.3.} \textit{Soit} $(\mathcal{M},W)$ \textit{une catégorie de modèles, associée à une catégorie à cofibrations. Alors sur le domaine} Diag \textit{formé des ensembles ordonnés finis, le prédérivateur} $\mathbf{D}=\mathbf{D}_{\mathcal{M},W}$ \textit{satisfait aux axiomes} Der (1, 2, 3, 4$^*$, 5$^*$, 6$^*$), \textit{et dualement. Donc si} $W$ \textit{est associé à la fois à une structure de cofibrations et de fibrations,} $\mathbf{D}$ \textit{est un dérivateur sur} Diag.

\msp{page 119}
\texttt{[Les six lignes qui suivent semblent hors contexte.]}

\textit{alors la flèche de changement de base}
\[ p^*f_!\longrightarrow f'_!q^* \]
\textit{entre foncteurs} $\mathrm{HOT}_W(X)\longrightarrow\mathrm{HOT}_W(S')$ \textit{est un isomorphisme.}

C'est pratiquement trivial, compte tenu, dans le cas $f$ $W$-lisse, que si $Z$ sur $X$ est $W$-lisse sur $X$, il l'est sur $S$ ; et dans le cas $p$ $W$-propre, en utilisant le dernier énoncé (``oubli'') dans \circled{6}, p. 199.

7.8. Je termine en mettant en vedette ce qui, à présent, me paraît la substance principale de tous ces ``rappels'' et comme la quintessence de la théorie des localiseurs fondamentaux
% margin: three vertical rules printed in the left margin, beside the footnotes and down to the foot of the page
% foot: Version 18 mars 2019
\printedpage{56}
% head: A. Grothendieck : Dérivateurs, Ch. XVI, Édition des Universités de Montpellier II et Paris VII
développée jusqu'à présent :
\[ \left\{ \begin{array}{rcl}
\text{Loc (3 bis)} & \Longrightarrow & \text{Der (4$^*$, 5$^*$)}\\[1ex]
\text{Loc (3 bis, 6 a)} & \Longrightarrow & \text{Der (1, 2, 3, 4$^*$, 5$^*$, 6$^*$)}\\[1ex]
\text{Loc (6$'$ a, 8, 9)} & \Longrightarrow & \text{tous les axiomes des dérivateurs (+ Loc (1 ter)).}
\end{array} \right. \]

Je considère les axiomes Loc (3 bis), Loc (8), Loc (9) comme des axiomes ``anodins'' - tous les localiseurs qu'on définit par voie cohomologique ou homotopique ordinaire, y satisfont. Il reste l'axiome Loc (6 a) ou Loc (6$'$ a) des $W$-cofibrations (forme faible),
\msp{page 120}
qui visiblement joue un rôle crucial dans la théorie de $\mathrm{HOT}_W$. Par contre, les axiomes Loc (4) (décomposition), Loc (5) (connexité) et Loc (7) et ses variantes (et bien sûr aussi Loc (6 b)) n'interviennent absolument pas ! C'est particulièrement impressionnant, en ce qui concerne l'axiome des $W$-fibrations Loc (7) et ses variantes (toutes équivalentes à Loc (7) en présence de Loc (8), à la seule exception de Loc (7 F)). Comme on a vu, Loc (7) n'est nullement impliqué par Loc (6 a, 8, 9) ($\mathbf{W}_0$ n'y satisfait pas), tout au plus se pourrait-il que Loc (7 F) le soit. Il en est de même de l'axiome de connexité Loc (5), comme on voit sur l'exemple de $W=\mathrm{Fl}(\mathrm{Cat})$, qui satisfait à tous les axiomes vus jusqu'à présent, à la seule exception
\msp{page 121}
de Loc (5). Quant à la question de l'indépendance de l'axiome Loc (4) par rapport à Loc (6 a, 8, 9), nous allons examiner ce genre de questions au paragraphe suivant.

L'ensemble d'axiomes Loc (6$'$ a, 8, 9) me paraît à présent si important, que j'ai envie de lui donner un nom, et d'appeler \textit{localiseur} (fondamental) \textit{spécial} dans Cat \footnote[57]{J'hésite si je ne dois pas rajouter l'axiome Loc (3 bis), qui assure que $W=W_{\mathbf{D}}$, $\mathbf{D}=\mathrm{HOT}_W$.}, un localiseur fondamental qui y satisfait. Ce sont donc aussi ceux pour lesquels nous savons que $\mathrm{HOT}_W$ est un vrai \textit{dérivateur.} Si on ne peut écarter comme ``pathologique'' une structure aussi riche et délicate qu'un dérivateur, il s'ensuit aussitôt que l'on ne peut non plus écarter les localiseurs spéciaux dans Cat (même si on affecte de regarder de haut les
\msp{page 122}
localiseurs fondamentaux généraux !). En effet, $W$ se récupère de façon unique à l'aide de $\mathbf{D}=\mathrm{HOT}_W$, ou de tout dérivateur équivalent à celui-ci, comme l'ensemble des $\mathbf{D}$-équivalences dans Cat \footnote[58]{Mais ceci n'est prouvé que moyennant Loc (3 bis).} : c'est le localiseur fondamental (correspondant à une notion d'équivalence ``universelle'', \textit{i.e.} dans la catégorie Cat de tous les ``espaces'' élémentaires) canoniquement associé au dérivateur. Ce n'est donc nullement un localiseur de référence fixe assez fin, tel $\mathbf{W}_\infty$, qui est apte à jouer un rôle !

Cela nous amène à nouveau à une question du genre déjà soulevé : si $W\subset\mathbf{W}_\infty$ est le plus petit localiseur qui satisfait à Loc (6$'$ a, 8) (ou Loc (3 bis, 6$'$ a, 8)), peut-on dire a) que $W$ est même \textit{spécial}, \textit{i.e.}
\msp{page 123}
satisfait également Loc (9) (accessibilité) ?, et de plus b) a-t-on $W\neq\mathbf{W}_\infty$ ? Je suis persuadé que c'est le cas (il y a moins de mérite pour ce $W$-là d'être différent de $\mathbf{W}_\infty$,
% foot: Version 18 mars 2019
\printedpage{57}
% head: A. Grothendieck : Dérivateurs, Ch. XVI, Édition des Universités de Montpellier II et Paris VII
que pour le $\tilde{\mathbf{W}}_\infty$ envisagé dans 6.4 (p. 87)).

7.9. Oubli : la question 1 bis (p. 98), sur l'existence d'une structure de Quillen close sur Cat, correspondant à $W$. Le seul résultat que je connaisse dans ce sens est celui de \textsc{Thomason}, construisant explicitement une telle structure de Quillen (de plus propre), dans le cas où $W=\mathbf{W}_\infty$. Je n'ai pas vraiment regardé la démonstration de \textsc{Thomason}, et ignore si celle-ci peut s'étendre à des localiseurs fondamentaux plus généraux, par exemple à ceux satisfaisant Loc (6$'$ a, 8, 9). S'il en était ainsi, cela prouverait d'ailleurs que ces localiseurs
\msp{page 124}
ont assez de cofibrations et assez de fibrations, donc qu'ils satisfont Loc (6 a) et Loc (7 F). (Pour ce qui est de Loc (6 a), il semble qu'on puisse le prouver directement très simplement, \textit{cf.} le paragraphe suivant. Mais ici encore, il faudrait regarder quelques exemples (équivalences homologiques ou cohomologiques, $\mathbf{W}_n$, \textit{i.e.} les $n$-équivalences d'homotopie), et regarder déjà si Loc (7 F) est satisfait.)

Il ne me semble pas que la théorie de \textsc{Thomason} ait été jusqu'à présent \texttt{[mots illisibles]} être très utile, techniquement parlant. Psychologiquement, je crois qu'il a été important, pour donner confiance que Cat est vraiment une excellente catégorie de modèles pour Hot. J'ai été vraiment enchanté et stimulé d'apprendre ce théorème. Mais la notion d'application de Dwyer (que j'ai apprise en 1983 dans l'article de \textsc{Thomason}) a été, techniquement, beaucoup plus importante pour moi : c'était exactement ce qui manquait à ma panoplie de notions et d'intuitions ``homotopiques'' dans Cat \dots
% foot: Version 18 mars 2019

\end{document}
